% ============================================================
% ch18.tex —— 第 18 章 四维 N = 2 规范理论（Spectral Networks, LNM 2386）
% 覆盖：印刷 pp.251–268（PDF p291–p308）
% 编号公式唯一真值：work/pages/pNNN.png（200 dpi）
% ============================================================

% [原文 p.251]
\chapter{四维 \texorpdfstring{$\mathcal N=2$}{N=2} 规范理论}

\ChapterAbstract{我们研究具有 $\mathcal N=2$ 超对称的四维量子场论。分析从这些理论的运动学结构开始，包括超多重态、拉格朗日量与经典模空间，随后讨论量子 $\mathcal N=2$ 规范动力学。}

四维 $\mathcal N=2$ 量子场论（QFT）以式 (17.1) 的 $\mathcal N=2$ 超庞加莱代数连同内禀对称性作为其无穷小对称代数。$\mathcal N=2$ 超庞加莱代数自然地包含 $\mathcal N=1$ 超庞加莱代数作为子代数，这使 $4d$ $\mathcal N=2$ QFT 成为 $4d$ $\mathcal N=1$ QFT 的一个更特殊的子集。与一般 $\mathcal N=1$ 情形的主要差别在于存在八个超荷 $Q^I_\alpha$ 与 $\bar Q^I_{\dot\alpha}$（$I=1,2$）。此外，四维 $\mathcal N=2$ 超对称容许一个复中心荷 $Z$，它出现在下述反对易关系中：
\[
\{Q^I_\alpha,Q^J_\beta\}=2\sqrt{2}\epsilon_{\alpha\beta}\epsilon^{IJ}Z
\qquad\text{and}\qquad
\{\bar Q^I_{\dot\alpha},\bar Q^J_{\dot\beta}\}=2\sqrt{2}\epsilon_{\alpha\beta}\epsilon^{IJ}Z^*\,.
\]
我们还在 $\mathcal N=2$ 理论的定义中要求它们具有一个整体 $\SU(2)_R$ $R$ 对称性，它按下式转动这些超荷：
\[
\text{for all }\alpha=1,2,\qquad
\begin{pmatrix} a & b\\ c & d\end{pmatrix}
\begin{pmatrix} Q^1_\alpha\\ Q^2_\alpha\end{pmatrix}
=
\begin{pmatrix} aQ^1_\alpha+bQ^2_\alpha\\ cQ^1_\alpha+dQ^2_\alpha\end{pmatrix},
\]
而对 $\bar Q^I_{\dot\alpha}$ 在共轭表示中有类似的作用。除 $\SU(2)_R$ 之外，$4d$ $\mathcal N=2$ 理论的 $R$ 对称性还可以包含一个 $\U(1)_R$ 因子，在其下 $Q^I_\alpha$ 的荷为 $-1$，而 $\bar Q^I_{\dot\alpha}$ 的荷为 $1$。然而，这一 $\U(1)_R$ 对称性并不严格属于 $\mathcal N=2$ 超庞加莱代数，它可能被明显破缺或者产生反常：只有 $\U(1)_R$ 的（可能是平凡的）子群才可能被实现。

% [Tutor Note: 本页第二个反对易关系里的 $\epsilon_{\alpha\beta}$ 用的是不带点的指标。按 Lorentz 协变性，$\{\bar Q^I_{\dot\alpha},\bar Q^J_{\dot\beta}\}$ 只能与带点的 $\epsilon_{\dot\alpha\dot\beta}$ 缩并（不带点指标无法与带点指标缩并），故原文此处应写作 $\epsilon_{\dot\alpha\dot\beta}$；这一处是排版笔误，本译文照原文译为 $\epsilon_{\alpha\beta}$，读者按 $\epsilon_{\dot\alpha\dot\beta}$ 理解即可。核对办法：把 $\mathcal N=1$ 时的同一关系与式 (17.1) 对照，$\bar Q$ 的指标全带点。]

\section{经典场论}

% [原文 p.252]
\subsection{超多重态}

在 $4d$ $\mathcal N=2$ 超对称理论中，规范与物质内容由只含自旋（或螺旋度）小于 $1$ 的洛伦兹粒子的 $\mathcal N=2$ 超多重态组成。这些超多重态通过由超荷导出的产生算符作用在 Clifford 真空上来构造（参看第 17.1.1 节）。此外，当相对于由超荷 $Q^1_\alpha$ 与 $\bar Q^1_{\dot\alpha}$ 生成的 $\mathcal N=1$ 子代数来看时，它们可以分解为 $\mathcal N=1$ 超多重态之和。注意只有阿贝尔子群
\[
\U(1)_J=\left\{\begin{pmatrix} \ee^{\ii\vartheta} & 0\\ 0 & \ee^{-\ii\vartheta}\end{pmatrix},\ \vartheta\in[0,2\pi)\right\}<\SU(2)_R
\]
保持这个特定的 $\mathcal N=1$ 子代数。

\subsubsection*{$\mathcal N=2$ 矢量多重态}

无质量 $\mathcal N=2$ 矢量多重态相对于任一 $\mathcal N=1$ 子代数分裂为一个无质量 $\mathcal N=1$ 矢量多重态 $V=(A_\mu,\lambda_\alpha)$ 与一个无质量 $\mathcal N=1$ 手征多重态 $\Phi=(\phi,\tilde\lambda_\alpha)$：
\[
(\mathbf{0},\mathbf{2}(\mathbf{1/2}),\mathbf{1})\oplus_{C\mathcal{PT}}(\mathbf{0},\mathbf{2}(\mathbf{-1/2}),-\mathbf{1})\,.
\]
若相应的 $\mathcal N=1$ 矢量超多重态是阿贝尔的，则对应的 $\mathcal N=2$ 超多重态也称为阿贝尔的。

$\SU(2)_R$ 对称性平凡地作用在 $A_\mu$ 与 $\phi$ 上，而把 $(\lambda_\alpha,\tilde\lambda_\alpha)$ 当作二重态作用。把 $\mathcal N=2$ 超多重态的庞加莱内容表示成菱形是方便的；对矢量超多重态，该菱形为：
\[
\begin{matrix}
& A_\mu &\\[2pt]
\lambda_\alpha & & \tilde\lambda_\alpha\,,\\[2pt]
& \phi &
\end{matrix}
\]
其中西南–东北（相应地，东南–西北）方向对应于 $Q^1_\alpha$ 与 $\bar Q^1_{\dot\alpha}$（相应地，$Q^2_\alpha$ 与 $\bar Q^2_{\dot\alpha}$）的作用，而 $\SU(2)_R$ 作用在这些行上。

正如 $\mathcal N=1$ 情形，给定带生成元 $T^a$ 的单紧李代数 $\g$，可以构造非阿贝尔 $\mathcal N=2$ $\g$-矢量超多重态：
\[
V=2g(V_a\oplus\Phi_a)T^a\,,
\]

% [原文 p.253]
其中 $g$ 是规范耦合。换言之，无质量 $\mathcal N=2$ $\g$-矢量超多重态由一个无质量 $\mathcal N=1$ $\g$-矢量超多重态与一个位于 $\g$ 的伴随表示中的无质量 $\mathcal N=1$ 手征多重态组成。

有质量 $\mathcal N=2$ 矢量超多重态分裂为一个有质量 $\mathcal N=1$ 矢量多重态与一个有质量 $\mathcal N=1$ 手征多重态，其庞加莱内容为：
\[
(\mathbf{-1},\mathbf{4}(\mathbf{-1/2}),\mathbf{6}(\mathbf{0}),\mathbf{4}(\mathbf{1/2}),\mathbf{1})\,.
\]
与 $\mathcal N=2$ 超希格斯机制相一致，这对应于一个无质量 $\mathcal N=2$ 矢量超多重态与一个无质量 $\mathcal N=2$ 超多重态之和，后者将在稍后描述。有质量 $\mathcal N=2$ 矢量超多重态可以是（半）BPS 的，此时其庞加莱内容为 $(\mathbf{-1},\mathbf{2}(\mathbf{-1/2}),\mathbf{2}(\mathbf{0}),\mathbf{2}(\mathbf{1/2}),\mathbf{1})$。

\subsubsection*{超多重态}

在本段及下文，我们常把（超）场 $\Phi$ 的厄米伴随记作 $\Phi^\dagger$ 而不是 $\bar\Phi$，以减轻记号负担。无质量 $\mathcal N=2$ 超多重态分裂为一个无质量 $\mathcal N=1$ 手征多重态\footnote{我们遵循的约定是：超多重态中的手征多重态记作 $Q$（Quark 的首字母）。它们表示 $\mathcal N=1$ 手征超场，而不是超荷——这一区别总能从上下文看清。} $Q=(q,\psi_\alpha)$ 与一个无质量 $\mathcal N=1$ 反手征多重态 $\tilde Q^\dagger=(\tilde q^\dagger,\tilde\psi^\dagger_{\dot\alpha})$ 之和。为了属于同一个 $\mathcal N=2$ 超多重态，手征多重态 $Q$ 与 $\tilde Q$ 必须具有相反的荷。于是，无质量超多重态的庞加莱内容为：
\begin{equation}
(\mathbf{-1/2},\mathbf{2}(\mathbf{0}),\mathbf{1/2})\oplus_{C\mathcal{PT}}(\mathbf{1/2},\mathbf{2}(\mathbf{0}),-\mathbf{1/2})\,, \tag{18.1}
\end{equation}
这可以改写成菱形：
\[
\begin{matrix}
& \psi_\alpha &\\[2pt]
q & & \tilde q^\dagger\\[2pt]
& \tilde\psi^\dagger_{\dot\alpha} &
\end{matrix}
\]
必须加上 $C\mathcal{PT}$ 共轭，源于这样的限制：$\SU(2)_R$ 必须把超多重态的标量分量当作二重态作用，而对出现在式 (18.1) 的两个副本之一的两个复标量来说，这永远不可能成立。然而，对于规范群 $G$ 的四元数型表示\footnote{群 $G$ 的复表示 $R$ 称为\emph{四元数型的}，若它容许一个反线性 $G$-等变映射 $j:R\to R$ 使 $j^2=-\mathrm{Id}_R$，等价地，容许一个反对称不变张量；称为\emph{实的}，若它容许一个反线性 $G$-等变映射 $j:R\to R$ 使 $j^2=\mathrm{Id}_R$，等价地，容许一个对称不变张量。不可约表示不可能同时是实的与四元数型的。}中的超多重态，存在把自由度数目减半的办法。

% [Tutor Note: 「四元数型对应 $\SO(2n)$、实型对应 $\mathrm{USp}(2n)$」这个对应与许多教材的直觉相反，容易记反，值得核对：以 $\mathcal N=2$ $\su(2)$ SYM 加 $n$ 个基础超多重态为例，$\su(2)$ 的基础表示是四元数型的（$\epsilon_{ab}$ 反对称不变），此时期望的味对称性是 $\SO(2n)$——正是 $\mathcal N=2$ $\su(2)$ SYM 加 $4$ 个味道时出现的 $\SO(8)$，也是该理论成为超共形场论的著名特例。反过来，伴随表示是实型的（$\SU(2)$ 的伴随）给 $\mathrm{USp}(2n)$，$n=1$ 时 $\mathrm{USp}(2)=\SU(2)$ 与 $\mathcal N=2^*$ 理论的 $\SU(2)$ 味对称性相符。两条判据分别是：$R$ 上不变张量对称则在 $2n$ 个副本上保持对称型不变形式（$\SO$），反对称则保持辛形式（$\mathrm{USp}$）。]

% [原文 p.254]
一般地，对在规范群 $G$ 的不可约复表示 $R$ 中变换的 $n$ 个无质量超多重态 $Q_i\oplus\tilde Q^{i\dagger}$，其味对称性是 $\SU(n)$。然而，若 $R$ 同构于其共轭表示 $\bar R$，则味对称性会增强。确切地说，当 $R$ 是四元数型时味对称性是 $\SO(2n)$，当 $R$ 是实型时味对称性是 $\mathrm{USp}(2n)$。当 $Q\oplus\tilde Q^\dagger$ 位于带反对称不变张量 $\epsilon_{ab}$ 的四元数型表示 $R$ 中时，可以施加：
\begin{equation}
\tilde Q_a=\epsilon_{ab}Q^b\,. \tag{18.2}
\end{equation}

\begin{definition}[18.1]
群 $G$ 的四元数型表示 $R$ 中的\emph{半超多重态}是满足式 (18.2) 约束的超多重态 $Q\oplus\tilde Q^\dagger$。
\end{definition}

有质量 $\mathcal N=2$ 超多重态必须是（半）BPS 的，才能只包含自旋小于 $1$ 的粒子，即它们必须满足 $m\ge\sqrt{2}|Z|$。半 BPS 超多重态的庞加莱内容为：
\[
(\mathbf{2}(\mathbf{-1/2}),\mathbf{4}(\mathbf{0}),\mathbf{2}(\mathbf{1/2}))\,,
\]
其中倍数 $2$ 同样由 $\SU(2)_R$ 不变性所要求。

\subsection{拉格朗日量}

我们主要关心 $\mathcal N=2$ 规范理论的低能行为。容许拉格朗日量描述的理论更易于分析，尽管并非所有 $\mathcal N=2$ 理论都允许这种描述。在拉格朗日型 $\mathcal N=2$ 规范理论中，可重整理论适合作为 UV 理论，而不可重整理论通常作为低能有效描述出现。

我们将用 $\mathcal N=1$ 超空间形式书写所有 $\mathcal N=2$ 拉格朗日量，从而按第 17.2.2 节由 $\mathcal N=1$ 手征超场与矢量超场来构造它们，参考的是由超荷 $Q^1_\alpha$ 与 $\bar Q^1_{\dot\alpha}$ 生成的 $\mathcal N=1$ 子代数。为确保理论保持 $\mathcal N=2$ 超对称，必须再施加 $\SU(2)_R$ 不变性这一附加条件。事实证明，$\SU(2)_R$ 对称性这一要求对于 $\mathcal N=1$ 拉格朗日量同时也具有 $\mathcal N=2$ 超对称而言既是必要的也是充分的。

\subsubsection*{可重整的 $\mathcal N=2$ 拉格朗日量}

\begin{proposition}[18.2]
具有单紧规范代数 $\g$ 的纯 $\mathcal N=2$ SYM 的拉格朗日量涉及一个 $\g$-非阿贝尔 $\mathcal N=1$ 矢量超场 $V$（其超对称场强为 $W_\alpha$）以及一个位于 $\g$ 的伴随表示中的 $\mathcal N=1$ 无质量手征超场 $\Phi$。它可写成：
\end{proposition}

% [原文 p.255]
\begin{equation}
\begin{aligned}
\mathcal L^{\mathcal N=2}_{\mathrm{SYM}}
&=\frac{1}{16\pi}\operatorname{Im}\left(\tau\int\de^2\theta\,\tr W^\alpha W_\alpha\right)
+2\tr\int\de^2\theta\de^2\bar\theta\,\bar\Phi \ee^{2g[V,\cdot]}\Phi\\
&=\tr\Big[-\frac{1}{2}F_{\mu\nu}F^{\mu\nu}-2\ii\lambda\sigma^\mu D_\mu\bar\lambda
-2\ii\tilde\lambda\sigma^\mu D_\mu\bar{\tilde\lambda}
+2\overline{D_\mu\phi}D^\mu\phi+\frac{\Theta g^2}{16\pi^2}F_{\mu\nu}\star F^{\mu\nu}\\
&\qquad+D^2+2\bar FF+2\ii\sqrt{2}g\bar\phi\{\lambda,\tilde\lambda\}
-2\ii\sqrt{2}g\{\bar{\tilde\lambda},\bar\lambda\}\phi+2gD[\phi,\bar\phi]\Big]\,,
\end{aligned}
\tag{18.3}
\end{equation}
其中记号见式 (17.18)、式 (17.19) 与式 (17.23)。

$\mathcal N=2$ Maxwell 理论由类似的拉格朗日量定义，只是 $V$ 为阿贝尔 $\mathcal N=1$ 矢量多重态而 $\Phi$ 为无质量中性手征超场。因此拉格朗日量的结构也相应改变：手征动能项中的 $\ee^{2g[V,\cdot]}$ 项消失，耦合 $\tau$ 变为式 (17.16) 定义的 $\tau_M$，而归一化按式 (17.15) 改变 $\ee^2/8\pi$。

在式 (18.3) 中，$\tr$ 如同式 (17.18)，应理解为 $\g$ 上的一个 Killing 型，相对于它长根的平方长度为二，且 $\g$ 的生成元 $T^a$ 在此表示中的归一化使得
\begin{equation}
\tr\left(T^aT^b\right)=\delta^{ab}/2\,. \tag{18.4}
\end{equation}
$\g$ 上非退化的度量 $\tr$ 诱导了 $\g$ 与其对偶之间的同构，对偶基为 $T_a=2\tr(T_a\,\cdot)\in\g^*$。这解释了式 (18.3) 中第二个超空间积分前那个因子 $2\tr$。分量展开中的交换子与反交换子来自 $\Phi$ 在伴随表示中变换这一事实：
\[
\begin{aligned}
\bar\phi\lambda\psi&=\bar\phi^c\lambda_a(T^a)^b{}_c\psi_b=\bar\phi_c\lambda_a(T^a)^{bc}\psi_b\\
&=2\bar\phi_c\tr\left(T^c[T^a,T^b]\right)\psi^b=2\tr\bar\phi\{\lambda,\psi\}\,.
\end{aligned}
\]
此外，若把这个拉格朗日量当作 $\mathcal N=1$ 拉格朗日量来处理，各项的相对归一化本可以是任意的。然而 $\mathcal N=2$ 超对称要求 $\SU(2)_R$ 不变性，这规定了 $\lambda$ 与 $\tilde\lambda$ 的动能项必须有匹配的归一化。这一对称性约束唯一地固定了 $\mathcal N=2$ 拉格朗日量中动能项的相对归一化。另外，由于 $\SU(2)_R$ 不变性，不能为 $\Phi$ 引入任何超势项。

在 $\mathcal N=2$ 理论中，物质场由位于 $\g$ 的某个表示 $R$ 中的超多重态 $Q_i\oplus\tilde Q^{i\dagger}$ 组成（$i$ 是味指标）。与前面一样，$V\oplus\Phi$ 是对应于 $\g$ 的 $\mathcal N=2$ 矢量超场。

\begin{proposition}[18.3]
此类超多重态的最一般的可重整 $\mathcal N=2$ 拉格朗日密度可由式 (17.23) 得到，它可写成：
\end{proposition}

% [原文 p.256]
\begin{equation}
\begin{aligned}
\mathcal L&=\int\de^2\theta\de^2\bar\theta\left(Q^{i\dagger}\ee^{2gV_R}Q_i+\tilde Q_i^\dagger\ee^{-2gV_R}\tilde Q^i\right)
+\sqrt{2}g\left(\int\de^2\theta\,\tilde Q^i\Phi_RQ_i+c.c.\right)\\
&\qquad+\left(\int\de^2\theta\,\mu^i{}_j\tilde Q^jQ_i+c.c.\right).
\end{aligned}
\tag{18.5}
\end{equation}
$\SU(2)_R$ 对称性下的不变性对拉格朗日量的结构施加了若干重要约束，包括各项的相对归一化、质量项的形式以及关于质量矩阵 $\mu$ 的特定条件。一个关键结果是：$\SU(2)_R$ 不变性要求质量矩阵 $\mu$ 满足条件
\begin{equation}
[\mu,\mu^\dagger]=0\,. \tag{18.6}
\end{equation}
这意味着质量矩阵 $\mu$ 必须是正规的，即它可以由一个味旋转对角化。于是这给出如下形式的对角质量项：
\[
\mu\tilde Q_aQ^a\,,
\]
其中 $a$ 是规范指标（在式 (18.5) 中被省略了）。对满足条件 $\tilde Q_a=\epsilon_{ab}Q^b$ 的半超多重态 $Q$，质量项必须为零：$\mu=0$。这意味着半超多重态本身不可能具有质量。此外，如第 17.1.1 节所讨论，$\mathcal N=2$ 理论中的有质量超多重态只包含自旋小于 $1$ 的粒子，因而必然是 BPS 态。最后，拉格朗日量中任何三次超势项都必须为零。这是因为 $\SU(2)_R$ 没有三次不变量，从而此类项与 $\SU(2)_R$ 不变性不相容，尽管可重整性考虑是允许它们的。

\subsubsection*{不可重整的 $\mathcal N=2$ 拉格朗日量}

若放弃 $\mathcal N=1$ 拉格朗日量的可重整性要求，一般会得到一个 $\mathcal N=1$ 规范非线性西格玛模型，它描述阿贝尔与非阿贝尔 $\mathcal N=1$ 矢量超场 $V=2gV_aT^a$ 以及在规范群下带荷的手征超场 $\Phi^i$，其拉格朗日量形如
\[
\begin{aligned}
\mathcal L&=\frac{1}{16\pi}\operatorname{Im}\left[\int\de^2\theta\,\tau_{ab}(\Phi)W^{\alpha,a}W^b_\alpha\right]
+\int\de^2\theta\de^2\bar\theta\,K(\Phi^i,(\bar\Phi \ee^{2gV})_i)\\
&\qquad+\int\de^2\theta\,W(\Phi^i)+\int\de^2\bar\theta\,\bar W(\bar\Phi^i)\,.
\end{aligned}
\]
规范耦合 $\tau$ 被提升为手征超场 $\Phi^i$ 的解析函数，因此它本身就是一个手征超场。若拉格朗日量在 $\mathcal N=2$ 超对称下不变，则部分手征超场 $\Phi^i$ 属于 $\mathcal N=2$ 矢量超场，而其余的手征超场配对成超多重态。$\SU(2)_R$ 不变性约束了规范动能项与 $\mathcal N=2$ 矢量多重态中手征超场动能项的相对归一化，

% [原文 p.257]
因为规范微子 $\lambda$ 与伴随手征场中的 Weyl 费米子 $\tilde\lambda$ 的动能项必须匹配。

在下文中，我们将主要关注作为可重整 $\mathcal N=2$ 理论模空间库仑分支上的低能有效理论而出现的不可重整 $\mathcal N=2$ 拉格朗日量（库仑分支即将定义）。这些低能理论按定义是无能隙的，并且只存在于阿贝尔库仑相中。

\begin{proposition}[18.4]
此类理论一般可以由只涉及阿贝尔规范场与无质量规范中性物质场的（不可重整）拉格朗日西格玛模型来描述。
\end{proposition}

\begin{proof}
有质量场按定义不贡献于低能有效理论。至于无质量带荷物质场，它们所带荷的规范群的阿贝尔规范耦合在红外（IR）趋于零，致使该规范群以及相关的带荷物质场退耦。此外，无质量带荷场不能参数化模空间，因为任何非零的真空期望值（VEV）都会自发破缺它们所带荷的规范群。因此，这些场以及那些阿贝尔规范群都不出现在低能有效拉格朗日量中。于是我们只剩下由阿贝尔规范场与在规范群下不带荷的无质量物质场构成的 $\mathcal N=2$ 拉格朗日西格玛模型。
\end{proof}

我们从仅由 $\mathcal N=2$ 阿贝尔矢量超场 $A_i=V_i\oplus\Phi_i$（$i=1,\dots,r$，其中 $r$ 是理论的秩；特别地，这里的 $i$ 不是味指标）构成的非线性西格玛模型开始讨论，遵循 [339, Sect. 2.4] 的论述。为后文作准备，我们此后把手征超场的标量分量记作 $a_i$ 而不是 $\phi_i$。要求 $\lambda_i$ 与 $\tilde\lambda_i$ 的动能项匹配，便得到下述关系：
\[
-\ii(\tau^{ij}-\bar\tau^{ij})=2\frac{\partial^2K}{\partial a_i\partial\bar a_j}\,.
\]
由于 $\tau^{ij}$ 关于 $a_i$ 是解析的，可知对任意 $i,j,k$ 都有
\[
-\ii\frac{\partial\tau^{ij}}{\partial a_k}=2\frac{\partial^3K}{\partial a_k\partial a_i\partial\bar a_j}\,,
\]
而它关于 $i,j,k$ 对称。因此局部地可以写成
\begin{equation}
\tau^{ij}=\int\de a_k\frac{\partial\tau^{ij}}{\partial a_k}
=\frac{\partial}{\partial a_i}\int\de a_k\,\tau^{kj}
=\frac{\partial^2}{\partial a_i\partial a_j}\int\de a_k\de a_l\,\tau^{kl}
=:\frac{\partial^2\mathcal F}{\partial a_i\partial a_j}\,, \tag{18.7}
\end{equation}
其中\emph{预势} $\mathcal F$ 是 $a_i$ 的全纯函数。事实上，$\mathcal F$ 可以看作 $A_i$ 的解析函数，于是超场

% [原文 p.258]
\[
T^{ij}=\frac{\partial^2\mathcal F}{\partial A_i\partial A_j}
\]
以 $\tau_{ij}$ 为标量分量。预势一般只是局部定义的，因此与其把它看成真正的全纯函数，不如把它理解为由 $a_i$ 参数化的流形上某个全纯线丛的全纯截面。用预势表示，凯勒势 $K$ 可写成
\begin{equation}
K=\operatorname{Im}\left(\frac{\partial\mathcal F}{\partial a_i}\bar a^i\right). \tag{18.8}
\end{equation}

\begin{definition}[18.5]
若一个凯勒流形的凯勒势可以写成上述形式，则称它为\emph{特殊凯勒流形} [137]。
\end{definition}

\begin{proposition}[18.6]
$\mathcal N=2$ 超对称蕴含：仅由矢量超场构成的 $\mathcal N=2$ 非线性西格玛模型的目标流形是特殊凯勒流形，其凯勒度量为
\[
\de s^2=\frac{\partial^2K}{\partial a^i\partial\bar a^j}\de a^i\de\bar a^j=\operatorname{Im}(\tau_{ij})\de a^i\de\bar a^j\,.
\]
\end{proposition}

% [Tutor Note: 式 (18.7) 的最后一步 $\tau^{ij}=\partial^2\mathcal F/\partial a_i\partial a_j$ 依赖于 $\partial\tau^{ij}/\partial a_k$ 关于 $i,j,k$ 对称——这正是上一式所给，于是混合偏导可交换、二次微分形式 $\tau_{ij}\de a^i\de a^j$ 局部恰当。可用式 (18.8) 反向核对，同时避开一个常见误解：对 $K=\operatorname{Im}\left(\partial\mathcal F/\partial a_i\,\bar a^i\right)$ 求 $\partial^2K/\partial a_j\partial\bar a_k=\operatorname{Im}(\tau_{kj})$，正是命题 18.6 中 $\de s^2$ 的第二个等号；而对 $a_j,a_k$ 的二阶导给出 $\operatorname{Im}\left(\bar a^i\,\partial^3\mathcal F/\partial a_i\partial a_j\partial a_k\right)$，与式 (18.7) 上面那个三阶导关系相印证，一般并不为零。因此「凯勒势的二阶全纯导数为零」是错的：凯勒度量的分量取自 $a$ 与 $\bar a$ 的混合二阶导，而二阶全纯导数是另一个（非零的）对象。另需注意预势 $\mathcal F$ 一般只在局部定义，整体上是某个全纯线丛的截面，因此 $\tau_{ij}$ 的整体单值性（模空间上的 monodromy）不能仅由这个局部公式读出。]

现在让我们把注意力转向仅涉及无质量、不带荷超多重态的不可重整 $\mathcal N=2$ 拉格朗日量，如 [339, Sect. 7.1] 所讨论。施加 $\SU(2)_R$ 不变性会对西格玛模型的目标流形施加约束，具体地要求它不仅容许一个复结构，而且容许一整族复结构——构成一个 $\mathbb{CP}^1$ 这么多的复结构。

\begin{proposition}[18.7]
任何 $4d$ $\mathcal N=2$ 无质量、不带荷超多重态西格玛模型的目标流形都是超凯勒流形。
\end{proposition}

关于超凯勒流形的一般讨论可参阅 [201]。

\begin{proof}
这来自下述事实：对每个超对称变换 $\delta_\alpha$ 以及每个手征超场 $\Phi=(\lambda,\phi)$，都有
\[
\delta^\dagger_{\dot\alpha}\delta_\alpha
\begin{pmatrix} \operatorname{Re}\phi\\ \operatorname{Im}\phi\end{pmatrix}
=\sigma^\mu_{\alpha\dot\alpha}\partial_\mu I
\begin{pmatrix} \operatorname{Re}\phi\\ \operatorname{Im}\phi\end{pmatrix},
\]
其中
\[
I=\begin{pmatrix} 0 & 1\\ -1 & 0\end{pmatrix};
\]
特别地，$I^2=-1$。设 $\delta^1_\alpha$ 与 $\delta^2_\alpha$ 是 $\mathcal N=2$ 超庞加莱代数定义中出现的两个独立超对称。那么，对每个满足 $|c_1|^2+|c_2|^2=1$ 的 $c=(c_1,c_2)\in\C^2$，$\delta^c_\alpha=c_1\delta^1_\alpha+c_2\delta^2_\alpha$ 也生成一个 $\mathcal N=1$ 子代数。此外，矩阵
% [原文 p.259]
\[
I^{(c)}=I_a n^a\,,\qquad\text{where}\qquad
n^a=\left(\bar c_1\ \ \bar c_2\right)\sigma^a
\begin{pmatrix} c_1\\ c_2\end{pmatrix},
\]
全都平方为 $-1$，并且 $I_1,I_2,I_3$ 满足四元数的关系。
\end{proof}

阿贝尔矢量多重态与无质量中性超多重态的一般 $\mathcal N=2$ 拉格朗日量为
\begin{equation}
\begin{aligned}
&\int\de^2\theta\,\frac{-\ii}{8\pi}\frac{\partial^2\mathcal F}{\partial\Phi_i\partial\Phi_j}W_{i,\alpha}W^\alpha_j+\text{c.c.}
+\ii\int\de^2\theta\de^2\bar\theta\left(\frac{\partial\mathcal F}{\partial\Phi^\dagger_i}\Phi_i-\Phi_i^\dagger\frac{\partial\mathcal F}{\partial\Phi_i}\right)\\
&\qquad+\int\de^2\theta\de^2\bar\theta\,K_h(Q^\dagger,Q)\,,
\end{aligned}
\tag{18.9}
\end{equation}
其中 $Q$ 表示超多重态中的所有手征多重态，而 $K_h$ 是\emph{超凯勒势}，它是 $Q^\dagger$ 与 $Q$ 的实解析函数。超多重态标量流形上的度量由该势导出。

\begin{proposition}[18.8]
西格玛模型目标空间（由 $\mathcal N=2$ 矢量超场与超多重态中标量的真空期望值（VEV）参数化）上的度量取局部乘积形式 $\mathcal M_V\times\mathcal M_H$。这里 $\mathcal M_V$ 是特殊凯勒流形，其度量由预势 $\mathcal F$ 导出；而 $\mathcal M_H$ 是超凯勒流形，其度量由 $K_h$ 确定。
\end{proposition}

\begin{proof}
这来自矢量多重态与超多重态的拉格朗日量项完全退耦这一事实，因为超多重态是不带荷的。
\end{proof}

\section{经典模空间}

如同一般的 $\mathcal N=1$ QFT（参看第 17.2 节），$\mathcal N=2$ 规范理论中的标量势表示为标量场的非负函数之和。$\mathcal N=2$ 规范理论中\emph{超对称真空的经典模空间}定义为由标量真空期望值（VEV）参数化的复仿射空间的子簇，在其上标量势为零，并模去规范对称性。虽然在一般 $\mathcal N=1$ QFT 中所有标量场都对应于手征超场的分量，但在 $\mathcal N=2$ 规范理论中可以一致地区分两类标量场：属于矢量多重态的和属于超多重态的。

对由式 (18.3) 与式 (18.5) 中拉格朗日量项刻画的可重整 $\mathcal N=2$ 理论，定义经典模空间的方程来自三个不同辅助场的变分。本节我们恢复先前的记号：$a$ 是规范指标，而 $i,j,\dots$ 是味指标。

% [原文 p.260]
首先，$\mathcal N=1$ 矢量超场中的辅助场 $D$ 出现在矢量超场、$\mathcal N=2$ 矢量中的伴随手征超场以及超多重态中手征超场的动能项的分量展开里。这导致下述 $D$ 项方程：
\begin{equation}
-D^a=[\phi,\bar\phi]^a+\left((q^i)^\dagger T_R^a q^i-\tilde q_i^\dagger T_R^a\tilde q_i\right)=0\,,
\qquad\text{for every }a=1,\dots,\dim\g . \tag{18.10}
\end{equation}
其次，$\mathcal N=2$ 矢量中伴随手征超场的辅助场的变分给出 $F$ 项方程：
\begin{equation}
F^a_{\mathrm{adj}}=\tilde q^iT_R^a q_i=0\,,\qquad\text{for every }a=1,\dots,\dim\g . \tag{18.11}
\end{equation}
最后，超多重态中所含 $\mathcal N=1$ 手征超场里辅助场的变分导致下述 $F$ 项方程：
\begin{equation}
F_{\mathrm{hyper},i}=\phi q_i+\mu_i^j q_j=0\,,\qquad
\widetilde F^i_{\mathrm{hyper}}=\tilde q^i\phi+\mu_j^i\tilde q^j=0\,,
\qquad\text{for every }i . \tag{18.12}
\end{equation}

\begin{definition}[18.9]
理论的\emph{经典模空间}定义为由场 $\phi$、$q_i$、$\tilde q_i$ 的 VEV（它们分别是 $\g$、$R$ 与 $\bar R$ 中的常向量）张成的仿射空间的子簇，满足条件
\[
\sum_{a=1}^{\dim\g}\left(\frac{g^2}{2}|D^a|^2+|F^a_{\mathrm{adj}}|^2\right)
+\sum_i\left(|F_{\mathrm{hyper},i}|^2+|\widetilde F^i_{\mathrm{hyper}}|^2\right)=0\,.
\]
\end{definition}

这些方程在 $\SU(2)_R$ 下不变，因此：

\begin{proposition}[18.10]
任何 $\mathcal N=2$ 理论中经典 $\mathcal N=1$ 超对称真空的模空间总是 $\mathcal N=2$ 超对称的。
\end{proposition}

定义任一 $\mathcal N=2$ 规范理论真空的经典模空间的方程可以重新整理成三个互不相同的集合（例如 $\mathcal N=2$ SQCD 的情形见 [26]）。在第一个集合中，方程只依赖于 $\mathcal N=1$ 矢量超场中包含的伴随标量：
\begin{equation}
[\phi^\dagger,\phi]=0\,. \tag{18.13}
\end{equation}
在第二个集合中，它们只依赖于超多重态中的标量：
\begin{equation}
\left(q^{i\dagger}T_R^a q_i-\tilde q^iT_R^a\tilde q_i^\dagger\right)=0\,,\qquad
\left(\tilde q^iT_R^a q_i\right)=0\,,\qquad
\text{for all }a=1,\dots,\dim\g\,, \tag{18.14}
\end{equation}
而在第三个集合中，它们同时依赖于两类标量场：
\begin{equation}
\begin{gathered}
\phi q_i+\mu_i^j q_j=0\,,\qquad \tilde q^i\phi+\mu_j^i\tilde q^j=0\,,\\
\phi^\dagger q_i+\mu_i^{\dagger j}q_j=0\,,\qquad \tilde q^i\phi^\dagger+\mu_j^{\dagger i}\tilde q^j=0\,.
\end{gathered}
\tag{18.15}
\end{equation}
把定义经典模空间的方程这样拆分，导致模空间中若干有趣子簇的定义，它们称为\emph{分支}。

% [Tutor Note: 本页指标的上下位置在原文中并不统一，照抄容易出错，特此标出三处：其一，式 (18.10) 第一项写成 $(q^i)^\dagger T^a_R q^i$（$i$ 为上指标），第二项写成 $\tilde q_i^\dagger T^a_R\tilde q_i$（$i$ 为下指标），而式 (18.14) 中同一对对象写成 $q^{i\dagger}T^a_R q_i$ 与 $\tilde q^iT^a_R\tilde q_i^\dagger$；由于 $q$ 与 $\tilde q$ 在共轭表示中成对出现，求和 $\sum_i$ 中「一上一下」缩并才是正确的，两种写法等价。其二，式 (18.12) 与 (18.15) 中 $\mu_i^{\dagger j}$、$\mu_j^{\dagger i}$ 的 $\dagger$ 是加在 $\mu$ 上的（即 $\mu$ 的厄米共轭的转置分量），不是加在指标上的。其三，同一批标量在第 18.2 节后文（定义 18.14 与 p.263）改用 $h$、$\tilde h$ 表示，见那里的批注。]

% [原文 p.261]
\subsubsection*{库仑分支}

\begin{definition}[18.11]
\emph{库仑分支}是模空间中满足 $q_i=\tilde q^i=0$ 的子簇。
\end{definition}

由定义可知，库仑分支完全由 $\mathcal N=2$ 矢量中伴随手征超场的标量分量 $\phi$ 的 VEV 参数化，且式 (18.13) 是唯一的约束。此外，规范等价的构型必须被等同。事实证明：

\begin{proposition}[18.12]
在规范变换下总可以假定 $\phi$ 属于 $\g$ 的一个固定的 Cartan 子代数 $\h$。
\end{proposition}

由于 $\h$ 在 $\g$ 中的正规化子模去 $\h$ 同构于 Weyl 群 $W(\g)$，$\mathcal N=2$ $\g$-规范理论的经典库仑分支是 $\h/W$。在 $C$ 的一般点处，规范代数破缺为 $\mathfrak{u}(1)^r\rtimes W$，其中 $r$ 是 $\g$ 的秩。这解释了「库仑分支」这一名称：经典地，$C$ 的任一一般点处的有效低能理论是一个（不可重整的）阿贝尔 $4d$ $\mathcal N=2$ 规范理论。第 18.1.2 节表明库仑分支上的低能物理完全由预势 $\mathcal F$ 决定，而 $C$ 是特殊凯勒流形。

可以用 $\phi_a$ 的 $\g$-不变多项式之集来参数化 $C$（例如见 [142]），这些多项式也称为\emph{库仑分支算符}。对任何有限维单李代数 $\g$，$W$-不变多项式环是有限生成的，且生成元之间没有关系。因此 $C$ 被库仑分支算符无关系地参数化。

\begin{example}[18.13]
考虑纯 $\mathcal N=2$ $\su(2)$ SYM，其模空间仅由 $C$ 构成。模去规范变换后总可以假定 $\mathcal N=2$ 矢量多重态中标量 $\phi$ 的 VEV 具有形式
\begin{equation}
\langle\phi\rangle=\alpha\sigma_3=\alpha
\begin{pmatrix} 1 & 0\\ 0 & -1\end{pmatrix}, \tag{18.16}
\end{equation}
\end{example}

其中 $\alpha$ 是复数，而对应于 $\alpha$ 与 $-\alpha$ 的真空被 Weyl 群 $W\simeq\Z_2$ 的规范作用等同。有 $C\simeq\C/\Z_2$，而库仑分支算符环为 $\C[u]$，其中 $u=\alpha^2=\tr\langle\phi\rangle^2$。当 $\alpha\ne 0$ 时规范代数自发破缺为 $\mathfrak{u}(1)\rtimes W$。

当 $\alpha\ne0$ 时，标量 $\phi$（其 VEV 为 $\langle\phi\rangle$）的动能项 $2\tr|D_\mu\phi|^2=2\tr|\partial_\mu\phi+[A_\mu,\phi]|^2$ 给出矢量场的质量项
\[
2\tr\left|[A_\mu,\langle\phi\rangle]\right|^2\,.
\]
由于
\[
\left[
\begin{pmatrix} A^0/2 & W^+/2\\ W^-/2 & -A^0/2\end{pmatrix}_\mu,
\begin{pmatrix} \alpha & 0\\ 0 & -\alpha\end{pmatrix}
\right]
=\begin{pmatrix} 0 & -\alpha W^+\\ \alpha W^- & 0\end{pmatrix}_\mu,
\]

% [原文 p.262]
$W^\pm$ 玻色子获得质量，规范群确实自发破缺为 $\mathfrak{u}(1)\rtimes W$。注意这一计算是经典的，只在可以忽略量子效应时成立。

经典预势为
\begin{equation}
\mathcal F_{\mathrm{cl}}=\frac{1}{2}\tau_{\mathrm{cl}}A^2\,, \tag{18.17}
\end{equation}
因为
\[
\tau_{\mathrm{cl}}=\frac{\partial^2\mathcal F}{\partial A\partial A}\,.
\]
预势可以（并且一般确实会）接受量子修正。

当加入基础超多重态 $(Q,\tilde Q)$ 时，库仑分支保持不变，因为它与这些超多重态无关。然而，式 (18.16) 所给的真空期望值（VEV）选择会通过拉格朗日量中的项 $\tilde Q\Phi Q$ 影响 $Q$ 与 $\tilde Q$ 的质量：
\begin{equation}
\sqrt{2}\tilde Q\langle\Phi\rangle Q+\mu\tilde Q Q
=\left(\tilde Q_1\ \tilde Q_2\right)
\begin{pmatrix} \sqrt{2}\alpha+\mu & 0\\ 0 & -\sqrt{2}\alpha+\mu\end{pmatrix}
\begin{pmatrix} Q_1\\ Q_2\end{pmatrix}.
\tag{18.18}
\end{equation}

\subsubsection*{希格斯分支}

为简单起见，设质量矩阵 $\mu$ 为零（事实上，只要有两个或更多质量相同的超多重态，希格斯分支就存在，例如见 [314]）。

\begin{definition}[18.14]
\emph{希格斯分支} $\mathcal H$ 由 $\phi=0$ 定义，因此仅由 $h$ 与 $\tilde h$ 的 VEV 参数化，并满足式 (18.14)，再模去规范变换。
\end{definition}

% [Tutor Note: 定义 18.14 改用 $h,\tilde h$ 而不再用 $q_i,\tilde q^i$，这是原文的记号切换，不是新的场：同一定义里说「满足式 (18.14)」，而式 (18.14) 完全是用 $q^{i\dagger}T_R^a q_i$、$\tilde q^iT_R^a q_i$ 写出的，故 $h\equiv q_i$、$\tilde h\equiv\tilde q^i$（对应关系按味指标一并等同）。核对办法：把 $h,\tilde h$ 代回式 (18.14) 应得到同一个方程；p.263「一般地 $\phi,h,\tilde h$ 都不为零」也正好对应「三组方程都起作用」中的 Higgs 组。此后直到本章末都用 $h,\tilde h$。]

对 $h$ 与 $\tilde h$ 的一般选择，规范群 $G$ 完全破缺，这一分支由此得名。经典地，可以预期 $\mathcal H$ 上的低能有效理论一般是仅含无质量中性超多重态的不可重整 $\mathcal N=2$ 西格玛模型，从而希格斯分支是超凯勒的。更多细节参阅 [26] 与 [339, Chap. 7]。

在规范代数为 $\g$、含 $k$ 个位于 $\g$ 的表示 $R$ 中的超多重态的 $\mathcal N=2$ 规范理论中，每个超多重态由两个复标量场组成，总共有 $2k\dim R$ 个复标量。可以证明经典希格斯分支是 $\C^{2k\dim R}$ 由复化规范代数 $\g_\C$ 作出的\emph{超凯勒商}，其复维数为 $2(k\dim R-\dim\g_\C)$。

\begin{example}[18.15]
在含 $F$ 个基础超多重态的 $\mathcal N=2$ $\su(2)$ SYM 中，经典希格斯分支的复维数为 $4F-2\dim\su(2)=4F-6$。特别地，当 $F=0,1$ 时 $\mathcal H$ 为空，而当 $F=2,3,4$ 时其复维数分别为 $2,6,10$。
\end{example}

% [原文 p.263]
\begin{example}[18.16]
考虑含两个超多重态 $(Q_i\oplus\tilde Q^{i\dagger})$（$i=1,2$）的 $\mathcal N=2$ Maxwell 理论，其荷为 $1$，如 [339, Sect. 7.3.1] 所述。希格斯分支由下述方程确定：
\[
q_1\tilde q^1+q_2\tilde q^2=0
\qquad\text{and}\qquad
|q_1|^2+|q_2|^2-|\tilde q^1|^2-|\tilde q^2|^2=0\,,
\]
\end{example}

模去 $\U(1)$ 规范转动。第一个方程可通过令 $(q_1,\tilde q^1)=(z,\tilde zt)$ 与 $(q_2,\tilde q^2)=(\tilde z,-zt)$ 来解。第二个方程施加条件 $|z|^2+|\tilde z|^2=|t|^2(|z|^2+|\tilde z|^2)$，从而 $|t|=1$。$\U(1)$ 对称性按 $(z,\tilde z,t)\to(\ee^{\ii\theta}z,\ee^{\ii\theta}\tilde z,\ee^{-2\ii\theta}t)$ 作用，使我们可固定 $t=1$。于是
\[
\mathcal H=\{(z,\tilde z)\in\C^2\}/\!\sim,
\]
其中 $(z,\tilde z)\sim(-z,-\tilde z)$，这意味着 $\mathcal H\simeq\C^2/\Z_2$ 并带有其自然的超凯勒度量。

\subsubsection*{混合分支}

在经典模空间中，一般地 $\phi$、$h$ 与 $\tilde h$ 都不为零，于是定义模空间的三组方程（包括式 (18.15)）都起作用。若低能自由度对应于一个带中性无质量物质的不可重整 $\mathcal N=2$ 阿贝尔规范非线性西格玛模型，则它可由式 (18.9) 所给形式的拉格朗日密度描述。正如命题 18.8 所强调的，目标流形上的度量必然是库仑分支上特殊凯勒度量与希格斯分支上超凯勒度量的局部乘积。希格斯分支与库仑分支只在模空间的原点相交，在那里所有 VEV 都为零。

\subsubsection*{分支的对称性定义}

存在许多已知但没有已知拉格朗日量描述的 $4d$ $\mathcal N=2$ 规范理论，上述模空间分析对它们并不适用。然而，即使没有拉格朗日量或对基本自由度的确切了解，也可以论证超对称真空模空间的存在，并基于对称性给出各分支的定义。

希格斯分支与库仑分支的区分尤其取决于整体 $\SU(2)_R$ 对称性因子的命运，而这又取决于哪一类标量场获得 VEV。$\mathcal N=2$ 矢量多重态中的标量 $\phi$ 在该对称性下不带荷，而超多重态中的 $q$ 与 $\tilde q^\dagger$ 按二重态变换。若超多重态场在某个真空中获得非零 VEV，则 $\SU(2)_R$ 自发破缺，从而可以通过把该真空映到另一个真空而非平凡地作用。相反，在库仑分支上 $\SU(2)_R$ 保持不破缺，并平凡地作用在真空上。因此，定义 18.11 可以推广为：

% [原文 p.264]
\begin{definition}[18.17]
\emph{库仑分支}是模空间中 $\SU(2)_R$ 保持不破缺的子簇。与拉格朗日情形相类比，任一 $\mathcal N=2$ 规范理论的\emph{秩}定义为其库仑分支的维数。
\end{definition}

希格斯分支可以用类似的方式对非拉格朗日理论定义。通过研究 $\SU(2)_R$ 之外的整体对称性（例如重子对称性），甚至可以定义希格斯分支的有趣子分支 [26]。按整体对称性的命运来定义模空间的分支，巩固了如下直觉：这些分支代表低能物理呈现一致性质的真空。

推广命题 18.12 之下所作的评注，人们相信任何 $4d$ $\mathcal N=2$ 理论（不必是拉格朗日型的）的库仑分支都是有限维复仿射空间对一个有限群 $G$ 的商，且库仑分支算符环是自由生成的 [339]。满足这一条件的有限群 $G$ 是\emph{复反射群}，这由 Chevalley–Shephard–Todd 定理确立，而 Weyl 群是其特殊情形。值得注意的是，存在一些颇为奇异的 $4d$ $\mathcal N=2$ 理论，其库仑分支是复有限维向量空间对一个非 Weyl 群的复反射群的商（例如 $\mathcal N=3$ 的 S 折叠 [160, 8, 161]）。

\section{量子规范动力学}

\subsection{重整化}

我们在第 17.3.1 节强调过，超对称 $4d$ QFT 中的微扰量子修正只能产生 D 项，而不能产生 F 项。因此，超势中任何耦合的重整化都可以通过手征超场的重整化来表达。此外，规范不变性保证项 $gV$（其中 $V$ 是任一矢量超场，$g$ 是对应的规范耦合）不发生重整化。在 $\mathcal N=2$ 拉格朗日规范理论中，属于同一 $\mathcal N=2$ 矢量超场的 $\mathcal N=1$ 矢量超场与手征超场的重整化必须相同；同样，同一超多重态中的手征与反手征超场的重整化也必须相同。由 $\SU(2)_R$ 不变性所要求的超势项 $\sqrt{2}g\tilde Q\Phi Q$ 进一步蕴含超多重态中的手征超场不重整化。因此：

\begin{proposition}[18.18]
希格斯分支上的超凯勒度量是树图精确的。
\end{proposition}

相反，库仑分支上的特殊凯勒度量可以、并且一般确实会接受量子修正。

由于 $\mathcal N=2$ SYM 理论是 $\mathcal N=1$ 理论的特殊情形，任一具有紧单规范群的 $\mathcal N=2$ SYM 理论的全纯规范耦合在微扰论中单圈精确：
\begin{equation}
\frac{\de}{\de\log\mu}\tau=-\frac{b_1}{2\pi\ii}+\text{non-perturbative}\,, \tag{18.19}
\end{equation}

% [原文 p.265]
其中 $b_1=3C(\mathrm{Adj})-n_\chi I(R_\chi)$ 是式 (17.27) 中的单圈重整化系数。（与标度无关的）全纯动力学标度为
\begin{equation}
\Lambda=\mu\exp\left(\frac{2\pi\ii\tau}{b_1}\right). \tag{18.20}
\end{equation}
在 $\mathcal N=2$ 理论中，每个 $\mathcal N=1$ 矢量超场都与一个伴随表示中的手征超场配对。因此 $b_1$ 的表达式化简为
\[
b_1=2C(\mathrm{Adj})-n_\chi I(R_\chi)\,,
\]
其中现在对 $\chi$ 的求和只遍历超多重态中的 $\mathcal N=1$ 手征超场。

\begin{example}[18.19]
在含 $F$ 个基础超多重态的 $\mathcal N=2$ $\su(N)$ SYM 中，即具有 $N$ 个颜色与 $F$ 个味道的 $\mathcal N=2$ SQCD，有 $b_1=2N-F$。特别地，当 $2F<N$ 时该理论渐近自由，而当 $2F>N$ 时红外自由。
\end{example}

\begin{example}[18.20]
在含一个伴随超多重态的 $\mathcal N=2$ $\su(N)$ SYM 中，有 $b_1=0$。当该超多重态有质量时，这些理论（常称为 $\su(N)$ $\mathcal N=2^*$ 理论）依赖于标度，因为超多重态在低于其质量的能量下标度退耦。相反，当伴随超多重态的质量为零时，理论不仅具有 $\mathcal N=2$ 而且事实上具有 $\mathcal N=4$ 超对称。这些理论称为 $\mathcal N=4$ $\su(N)$ SYM。它们是标度不变的 [316]，甚至是超共形的。
\end{example}

% [Tutor Note: 本章先后出现两个不同的 $b_1$ 公式，务必分开记。一般（$\mathcal N\ge1$）公式是 $b_1=3C(\mathrm{Adj})-\sum_\chi n_\chi I(R_\chi)$：规范微子贡献 $3C(\mathrm{Adj})$，每个手征多重态贡献 $-I(R_\chi)$。在 $\mathcal N=2$ 中每个 $\mathcal N=1$ 矢量超场都配了一个伴随手征超场，把规范微子与这个伴随手征场的贡献合并后得到 $b_1=2C(\mathrm{Adj})-\sum_{\chi\in\text{hyper}}n_\chi I(R_\chi)$，此时求和只遍历超多重态中的手征超场。可用两个例子核对（取 $C(\mathrm{Adj})=N$、$I(\text{基础})=1$、$I(\mathrm{Adj})=2N$ 的归一化）：纯 $\su(N)$ SYM 给 $b_1=2N$；加 $F$ 个基础超多重态给 $b_1=2N-F$（例 18.19），故 $2F<N$ 渐近自由；加一个伴随超多重态给 $b_1=2N-2N=0$（例 18.20），故 $\mathcal N=2^*$ 与 $\mathcal N=4$ 理论都不跑动。]

在第 17.3.4 节中，我们强调了 $\mathcal N=1$ 规范理论中全纯规范耦合与物理规范耦合的区别。由于两者之间的变换具有非解析性，一般 $\mathcal N=1$ 规范理论中的物理规范耦合在微扰论的各阶都接受修正。然而，由于 $\mathcal N=2$ 矢量多重态中存在伴随手征超场，这一现象在 $\mathcal N=2$ 理论中不会发生 [31]。

\begin{proposition}[18.21]
在 $\mathcal N=2$ 规范理论中，物理规范耦合与全纯规范耦合重合，从而在微扰论中单圈精确，并且当单圈系数为零时树图精确。
\end{proposition}

例如这意味着具有 $N$ 个颜色与 $2F$ 个味道的 $\mathcal N=2$ SQCD 中的规范耦合在微扰论中树图精确；此外，人们普遍相信非微扰量子修正也消失。这些理论构成另一类标准的 $\mathcal N=2$ SCFT [314, 146]。

\begin{remark}[18.22]
在渐近自由的拉格朗日型 $4d$ $\mathcal N=2$ 理论的库仑分支上，式 (18.19) 中规范耦合的重整化可以重写成对支配低能动力学的预势的修正，例如见 [316]。此外，这一重整化会传播到低能阿贝尔规范耦合。由于库仑分支上的特殊凯勒度量依赖于后者，可见库仑分支度量确实接受量子修正。
\end{remark}

% [原文 p.266]
\subsection{反常}

\subsubsection*{规范反常}

\begin{proposition}[18.23]
拉格朗日型 $4d$ $\mathcal N=2$ 规范理论中没有小的手征规范反常。
\end{proposition}

\begin{proof}
在拉格朗日型 $4d$ $\mathcal N=2$ 规范理论中，所有物质场或者属于 $\mathcal N=2$ 矢量超场中的伴随手征超场，或者属于超多重态中的手征超场。这些场不贡献于微扰手征规范反常。这是因为，第一，伴随表示是实的；第二，每个超多重态由两个位于规范群共轭表示中的手征超场组成，它们对反常的合并贡献相互抵消。
\end{proof}

然而，可以存在整体规范反常，例如当 $G=\mathrm{USp}(2N)$ 时的 Witten 反常。例如，含一个位于基础表示中的半超多重态的 $\mathcal N=2$ $\SU(2)$ SYM 作为量子理论是不自洽的 [363]。

\subsubsection*{阿贝尔整体对称性的反常}

手征 $\U(1)$ 整体对称性群可以产生反常，参看第 17.3.3 节。该反常为
\[
\mathcal A=\sum_i q_i I(R_i)\,,
\]
其中求和遍历在该手征对称性下荷为 $q_i$、且位于规范群表示 $R_i$ 中的 Weyl 费米子。

考虑纯 $\SU(N)$ $\mathcal N=2$ SYM 中的 $\U(1)_R$ $R$ 对称性（与 $\SU(2)_R$ 对称性不同），在其下场 $A_\mu,\lambda,\tilde\lambda,\phi$ 分别带荷 $0,1,1,2$。该对称性是反常的：
\[
\partial_\mu j^\mu_R=\frac{\mathcal A}{16\pi^2}F^{\mu\nu}_a\widetilde F^a_{\mu\nu}\,,
\]
其中 $\mathcal A=2I(\mathrm{adj})=2N$。参数为 $\alpha$ 的 $\U(1)_R$ 变换与 $\Theta$ 角的移动效果相同：$\Theta\to\Theta+2\alpha\mathcal A=\Theta+4N\alpha$。因此，只有 $\U(1)_R$ 的 $(\Z_{4N})_R$ 子群才是 QFT 的真正对称性，因为它对应于把 $\Theta$ 移动 $2\pi$ 的整数倍。此外，由于 $\phi$ 在 $(\Z_{4N})_R$ 下带荷 $2$，该对称性在库仑分支上自发破缺为 $\Z_2$。

与 $\mathcal N=1$ SQCD 不同，添加物质不一定恢复完整的 $R$ 对称性。由于 $R(\phi)=2$，超势项 $H_1\Phi H_2$ 迫使超多重态具有 $R$ 荷 $0$。结果，它们的 Weyl 分量带荷 $-1$，并贡献于 ABJ 反常。例如，具有 $N$ 个颜色与 $F$ 个基础味道的 $\mathcal N=2$ SQCD 中的量子 $R$ 对称性是 $\Z_{4N-2F}$。然而，当 $F=2N$ 时

% [原文 p.267]
$R$ 对称性是无反常的，这与该理论是 SCFT（其中 $\U(1)_R$ 因子是 $\mathcal N=2$ $4d$ 超共形代数的一部分）这一事实相一致。

\subsection{模空间}

考虑一个带无质量超多重态的渐近自由 $\mathcal N=2$ 规范理论。在远离模空间原点的超对称真空中，规范群在 UV 理论的弱耦合区域自发破缺，或者完全破缺，或者破缺到阿贝尔子群。此时经典分析给出理论在高于破缺标度的能标上真实动力学的一个良好近似。在库仑分支上，所得的低能理论通常是由阿贝尔 $\mathcal N=2$ 矢量超场与中性无质量超多重态构成的自由理论，它不随重整化群流跑动。因此模空间的经典描述在远离其原点处是一个良好近似。然而，当接近原点时，强耦合的非阿贝尔规范相互作用开始起作用，使经典分析失效。这一推理也适用于具有质量超多重态的理论（它们在低于其质量标度时退耦）以及 SCFT（当规范耦合小时它们是弱耦合的）。

我们在命题 18.8 中看到，任何 $4d$ $\mathcal N=2$ QFT 的经典模空间局部地是库仑分支与希格斯分支的乘积。这对量子模空间同样成立：若低能有效理论由 $\mathcal N=2$ 矢量超场与中性无质量超多重态构成，则量子模空间仍是特殊凯勒流形（库仑分支）与超凯勒流形（希格斯分支）的局部乘积 [26]。然而，由于量子修正，这些分支上的量子度量可能不同于其经典对应物。不过，有一个重要性质：

\begin{proposition}[18.24]
希格斯分支是经典精确的。
\end{proposition}

换言之，希格斯分支上的量子超凯勒度量与经典度量一致。这可以通过 $\mathcal N=2$ 全纯性来理解，其中 UV 耦合被提升为背景伪超场 [26]。要点在于希格斯分支的度量不依赖于 $\mathcal N=2$ 矢量多重态。由于规范耦合 $\tau$（或等价地，动力学标度 $\Lambda$）以及质量 $\mu_i^j$ 都可以解释为伪 $\mathcal N=2$ 矢量超场的 VEV，希格斯分支度量与 $\Lambda$ 及 $\mu_i^j$ 无关。特别地，若某个希格斯分支经典地存在，则它在完整的量子理论中也存在。

相反，库仑分支上的度量一般接受量子修正。尽管如此：

\begin{proposition}[18.25]
库仑分支不会被量子效应抬起。
\end{proposition}

在远离库仑分支原点的半经典区域，经典分析给出量子动力学的一个良好近似。平坦方向由低能阿贝尔

% [原文 p.268]
$\mathcal N=2$ 矢量超场中标量的 VEV 参数化，而唯一可能抬起库仑分支的机制会涉及这些标量获得质量。然而，这只可能通过希格斯机制发生，而由于低能理论只由阿贝尔矢量场与无质量中性物质构成，该阿贝尔规范群不可能被进一步希格斯化。因此，库仑分支在量子理论中存活下来，至少在远离原点处如此。通过解析延拓，这意味着即使 UV 的 $\mathcal N=2$ 矢量超场中伴随标量的 VEV 很小，库仑分支依然存在。因此，拉格朗日型 $\mathcal N=2$ 规范理论容许一个量子的库仑分支，它是特殊凯勒流形。其度量在远距离处渐近地接近经典度量。量子效应如何修改库仑分支度量将在第 20 章中详细探讨。

% [Tutor Note (全章汇总): 本章有若干处记号切换、指标位置不统一与一处笔误，阅读时留意：(1) 超多重态的标量在式 (18.2)—(18.15) 中记为 $q_i,\tilde q^i$，但在定义 18.14 与 p.263 正文中改用 $h,\tilde h$（原文如此），二者指同一批标量分量；(2) 式 (18.10) 与 (18.14) 的指标上下位置不统一，但求和 $\sum_i$ 中「一上一下」缩并即可，详见 p.260 的批注；(3) p.251 的 $\{\bar Q,\bar Q\}$ 反对易子应带带点指标 $\epsilon_{\dot\alpha\dot\beta}$，详见 p.251 的批注；(4) 命题 18.7 的证明末尾，原文把 $I^{(c)}=I_a n^a$ 这一段排在 p.259 页首、与 p.258 的证明断开，本译文按原顺序保留，读者应把它读作该证明的延续。]
